\subsubsection{Shape of the candidate curve}

Before matching the bounds, we verify that the six displayed branches have
the advertised shape.  The only nontransparent part is the algebraic branch
$\Rquad$; the remaining monotonicity follows directly from the reciprocal
formula and the mixed parametrization.

\begin{lemma}[Monotonicity of the deterministic curve]
\label{lem:opt:curve-monotonicity}
The function $\Rquad(u)$ is strictly decreasing on
$[\uDet{2},\uDet{3}]$.  The reciprocal and mixed branches are strictly
decreasing on their nonconstant ranges.  Consequently the six-branch
function in \cref{eq:opt:full-curve} attains its global maximum only at
$u=\uDet{2}$.
\end{lemma}

\begin{proof}
Implicit differentiation of $T_s(\Rquad(u))=0$, with $s=u-1$, gives
\begin{equation}
 \Rquad'(u)=
 -\frac{R^2+(3s^2+2s)R-3(s+1)^2}
 {2sR+s^3+s^2+1}.
 \label{eq:opt:Rquad-derivative}
\end{equation}
Using $T_s(R)=0$, its numerator factors as
\[
 \frac{2s^3+s^2-1}{s}
 \left(R-
 \frac{2s^3+3s^2-2}{2s^3+s^2-1}\right).
\]
On the relevant interval $s\ge \uDet{2}-1>2/3$, the prefactor is
positive, while the fraction in parentheses is at most $5/3$ because
$4s^3-4s^2+1>0$.  The sign certificate
\eqref{eq:opt:T-five-thirds} gives $R>5/3$, and hence
$\Rquad'(u)<0$.  The reciprocal branch is visibly decreasing.  In the
mixed parametrization, $m(c)$ is increasing and
$u(c)=1+2/c-m(c)$ is decreasing, so $1+c$ decreases with $u$.
The first nonconstant branch increases up to $u=\uDet{2}$, all subsequent
nonconstant branches decrease, and the last branch is constant.  This proves
the global-maximum assertion.
\end{proof}

\subsubsection{Exact joins and completion of the proof}

We next check that adjacent formulas agree at every transition point.  Once
this consistency is recorded, the theorem follows by pairing the upper and
lower result already proved for each interval.

\begin{lemma}[Exact joins]
\label{lem:opt:curve-joins}
The six branches in \cref{eq:opt:full-curve} agree at all five transition
points.  More explicitly, the first two branches both equal one at
$\uDet{1}$, and
\begin{align}
 \Rquad(\uDet{2})&=\uDet{2},
 \notag\\
 \Rquad(\uDet{3})&=1+\frac1{\sqrt{s_0}}
              =1.682327803828019\ldots,
 \notag\\
 1+\frac1{\sqrt{\uDet{4}-1}}&=\varphi
              =\Rmix(\uDet{4}),
 \notag\\
 \Rmix(\uDet{5})&=1+r=\RdetOT.
 \label{eq:opt:curve-joins}
\end{align}
\end{lemma}

\begin{proof}
The equality at $\uDet{1}=1$ is immediate.  Substitution of the defining
equations for $\uDet{2}$ and $s_0=\uDet{3}-1$ into $T_s$ gives the first
two displayed identities; uniqueness of the positive root identifies the
values with $\Rquad$.  Since
$\uDet{4}-1=\varphi+1=\varphi^2$, the reciprocal branch equals
$1+1/\varphi=\varphi$.  The two remaining equalities are exactly the endpoint
conditions of the mixed parametrization in
\cref{eq:opt:mixed-endpoints}.
\end{proof}

\begin{proof}[Proof of \cref{thm:optional-curve}]
For the upper bounds, \cref{lem:opt:raw-upper} gives the first two branches,
\cref{lem:opt:ute-endpoint-game} gives the algebraic branch, and
\cref{lem:opt:zero-prefix} gives the reciprocal branch.  The mixed branch
and the uniform plateau follow from
\cref{thm:opt:optional-middle-upper}.

For the matching lower bounds, \cref{prop:opt:binary-lower} supplies the
first four branches, \cref{prop:opt:mixed-lower} supplies the fifth, and
\cref{lem:opt:obligatory-transfer,thm:lower} supplies the plateau with
cap-free instances.  The exact joins are given by
\cref{lem:opt:curve-joins}, and the claimed shape follows from
\cref{lem:opt:curve-monotonicity}.  The finite upper bounds have the stated
$O_u(n)$ remainders, while
\cref{thm:opt:optional-middle-upper} gives one absolute $O(n)$ remainder on
the plateau.  This proves every assertion of the theorem.
\end{proof}

\begin{remark}[Representative mixed values]
For orientation, some values on the mixed branch are
\begin{center}
\begin{tabular}{@{}rrrr@{}}
\toprule
$u$ & $\RdetRO(u)$ & $m$ & $t$\\
\midrule
$3.800000$ & $1.598851588$ & $0.539725635$ & $1.138984621$\\
$4.000000$ & $1.583283633$ & $0.428863572$ & $1.330517832$\\
$4.071679486$ & $1.579204806$ & $0.381330539$ & $1.412236296$\\
$4.300000$ & $1.571692026$ & $0.198387083$ & $1.734651379$\\
\bottomrule
\end{tabular}
\end{center}
\end{remark}

\begin{remark}[Asymptotic scope]
\label{rem:opt:curve-audit}
All statements above concern the fixed-$u$ size-asymptotic ratio.  Below the
plateau, the upper bounds have additive error $O_u(n)$.  On
$u\in[\uDet{5},\infty)$ the same algorithm has one uniform $O(n)$ error.
The unbounded obligatory endpoint is not included in the statement of the
curve; it is derived in \cref{sec:obligatory}.
\end{remark}
